% GENERATED FILE. Edit canonical lessons, then run npm run generate:books.
% canonical-source-hash: fnv1a64:1861c6f60ff6e323
% canonical-lessons: MR-C157-cakhne, MR-C157-perne, MR-C157-vadhne, MR-C157-thakne, MR-C157-jagne

\chapter{To taste, To sow, To grow, To get tired, To stay awake}
\label{ch:mr-to-taste-to-sow-to-grow-to-get-tired-to-stay-awake}
\hlchaptermodality{157}

\begin{chapteropening}
\textbf{By the end of this chapter:} \emph{I can say to taste, to sow, to grow, to get tired and to stay awake.}

Complete the last lesson of chapter 157: i can say to taste, to sow, to grow, to get tired and to stay awake.
\end{chapteropening}

\section[cākhṇe]{\mr{चाखणे} (cākhṇe) — to taste}

\label{lesson:MR-C157-cakhne}



\begin{quote}
Before the new one: say the Marathi for to bite, then the Marathi for to swallow.
\end{quote}

\subsection*{You'll want to know: \mr{चाखणे}}
\textbf{\mr{चाखणे}} — \emph{cākhṇe} — \textquotedblleft{}to taste\textquotedblright{}.

\begin{grammarlens}[title={What you've built}]
One more word for this chapter.
\end{grammarlens}

\subsection*{Your turn}
\begin{itemize}
\raggedright
  \item \emph{Say it:} \emph{cākhṇe}
  \item \emph{Say it:} \emph{cākhṇe}, once more
  \item \emph{Say it:} say \emph{cākhṇe}
  \item \emph{Recall:} say the Marathi for to bite, then the Marathi for to swallow, then say \emph{cākhṇe} again
\end{itemize}

\begin{tcolorbox}[breakable,colback=teal!4,colframe=teal!35!black,title={Before you move on}]
What does \mr{चाखणे} mean? (\textquotedblleft{}To taste\textquotedblright{}.) Say it once more.
\end{tcolorbox}
\section[perṇe]{\mr{पेरणे} (perṇe) — to sow}

\label{lesson:MR-C157-perne}



\begin{quote}
Before the new one: say the Marathi for to swallow, then the Marathi for to taste.
\end{quote}

\subsection*{You'll want to know: \mr{पेरणे}}
\textbf{\mr{पेरणे}} — \emph{perṇe} — \textquotedblleft{}to sow\textquotedblright{}.

\begin{grammarlens}[title={What you've built}]
One more word for this chapter.
\end{grammarlens}

\subsection*{Your turn}
\begin{itemize}
\raggedright
  \item \emph{Say it:} \emph{perṇe}
  \item \emph{Say it:} \emph{perṇe}, once more
  \item \emph{Say it:} say \emph{perṇe}
  \item \emph{Recall:} say the Marathi for to swallow, then the Marathi for to taste, then say \emph{perṇe} again
\end{itemize}

\begin{tcolorbox}[breakable,colback=teal!4,colframe=teal!35!black,title={Before you move on}]
What does \mr{पेरणे} mean? (\textquotedblleft{}To sow\textquotedblright{}.) Say it once more.
\end{tcolorbox}
\section[vāḍhṇe]{\mr{वाढणे} (vāḍhṇe) — to grow}

\label{lesson:MR-C157-vadhne}



\begin{quote}
Before the new one: say the Marathi for to taste, then the Marathi for to sow.
\end{quote}

\subsection*{You'll want to know: \mr{वाढणे}}
\textbf{\mr{वाढणे}} — \emph{vāḍhṇe} — \textquotedblleft{}to grow\textquotedblright{}.

\begin{grammarlens}[title={What you've built}]
One more word for this chapter.
\end{grammarlens}

\subsection*{Your turn}
\begin{itemize}
\raggedright
  \item \emph{Say it:} \emph{vāḍhṇe}
  \item \emph{Say it:} \emph{vāḍhṇe}, once more
  \item \emph{Say it:} say \emph{vāḍhṇe}
  \item \emph{Recall:} say the Marathi for to taste, then the Marathi for to sow, then say \emph{vāḍhṇe} again
\end{itemize}

\begin{tcolorbox}[breakable,colback=teal!4,colframe=teal!35!black,title={Before you move on}]
What does \mr{वाढणे} mean? (\textquotedblleft{}To grow\textquotedblright{}.) Say it once more.
\end{tcolorbox}
\section[thakṇe]{\mr{थकणे} (thakṇe) — to get tired}

\label{lesson:MR-C157-thakne}



\begin{quote}
Before the new one: say the Marathi for to sow, then the Marathi for to grow.
\end{quote}

\subsection*{You'll want to know: \mr{थकणे}}
\textbf{\mr{थकणे}} — \emph{thakṇe} — \textquotedblleft{}to get tired\textquotedblright{}.

\begin{grammarlens}[title={What you've built}]
One more word for this chapter.
\end{grammarlens}

\subsection*{Your turn}
\begin{itemize}
\raggedright
  \item \emph{Say it:} \emph{thakṇe}
  \item \emph{Say it:} \emph{thakṇe}, once more
  \item \emph{Say it:} say \emph{thakṇe}
  \item \emph{Recall:} say the Marathi for to sow, then the Marathi for to grow, then say \emph{thakṇe} again
\end{itemize}

\begin{tcolorbox}[breakable,colback=teal!4,colframe=teal!35!black,title={Before you move on}]
What does \mr{थकणे} mean? (\textquotedblleft{}To get tired\textquotedblright{}.) Say it once more.
\end{tcolorbox}
\section[jāgṇe]{\mr{जागणे} (jāgṇe) — to stay awake}

\label{lesson:MR-C157-jagne}



\begin{quote}
Before the new one: say the Marathi for to grow, then the Marathi for to get tired.
\end{quote}

\subsection*{You'll want to know: \mr{जागणे}}
\textbf{\mr{जागणे}} — \emph{jāgṇe} — \textquotedblleft{}to stay awake\textquotedblright{}.

\begin{grammarlens}[title={What you've built}]
One more word for this chapter.
\end{grammarlens}

\subsection*{Your turn}
\begin{itemize}
\raggedright
  \item \emph{Say it:} \emph{jāgṇe}
  \item \emph{Say it:} \emph{jāgṇe}, once more
  \item \emph{Say it:} say \emph{jāgṇe}
  \item \emph{Recall:} say the Marathi for to grow, then the Marathi for to get tired, then say \emph{jāgṇe} again
\end{itemize}

\begin{tcolorbox}[breakable,colback=teal!4,colframe=teal!35!black,title={Before you move on}]
What does \mr{जागणे} mean? (\textquotedblleft{}To stay awake\textquotedblright{}.) Say it once more.
\end{tcolorbox}
